\appendix

\section{Model capability index}
\label{app:capability}

\citet{kim2026scaling} score models on a published Intelligence Index and report that a
\emph{task-grounded} capability metric explains more variance ($R^2 = 0.413$) than the generic one
($R^2 = 0.373$). We therefore use the task-grounded form throughout: $I$ is a model's mean
single-doctor chain-of-thought accuracy across the three benchmarks. Table~\ref{tab:capability}
gives the components. The three tiers span $I = 34.0$ to $59.2$, and their single-benchmark
baselines span 15.6\% to 91.6\% --- crossing the 45\% region from both sides, which is what makes
the capability-ceiling test possible at all.

\begin{table}[h]\centering\small
\caption{Capability index and its components (single-doctor CoT accuracy, \%).}
\label{tab:capability}
\setlength{\tabcolsep}{4pt}
\renewcommand{\arraystretch}{1.02}
\begin{tabular}{@{}llrrrrl@{}}
\toprule
Vendor & Model & MedXpertQA & MedAgentsBench & MedQA & $I$ (mean) & Setting \\
\midrule
\sectrow
\multicolumn{7}{@{}l}{\textit{OpenAI --- the capability ladder the grid is run on}} \\
OpenAI & \texttt{gpt-4.1-nano} & 15.6 & 20.4 & 66.0 & \textbf{34.0} & --- \\
\winrow
OpenAI & \texttt{gpt-5-nano}   & 32.0 & 32.4 & 88.0 & \textbf{50.8} & effort=low \\
\winrow
OpenAI & \texttt{gpt-5-mini}   & 41.6 & 44.4 & 91.6 & \textbf{59.2} & effort=low \\
\midrule
\sectrow
\multicolumn{7}{@{}l}{\textit{Google --- the second vendor, for the diversity ladder}} \\
\winrow
Google & \texttt{gemini-3.5-flash-lite} & 34.8 & 31.2 & 84.8 & \textbf{50.3} & --- \\
Google & \texttt{gemini-3.7-flash}      & 60.4 & 56.4 & 95.2 & \textbf{70.7} & --- \\
\bottomrule
\multicolumn{7}{@{}l}{\scriptsize shaded rows fall inside the 25--50\% collaboration window on at} \\
\multicolumn{7}{@{}l}{\scriptsize least one benchmark. \texttt{gemini-3.5-flash-lite} ($I{=}50.3$) and} \\
\multicolumn{7}{@{}l}{\scriptsize \texttt{gpt-5-nano} ($I{=}50.8$) are the capability-matched cross-vendor pair.} \\
\end{tabular}
\end{table}

\section{Domain complexity}
\label{app:complexity}

\subsection{Complexity metric construction}
Following \citet{kim2026scaling}, domain complexity $D \in [0,1]$ is the arithmetic mean of three
complementary measures: the \textbf{performance ceiling} $1 - p_{\max}$, where $p_{\max}$ is the
highest performance any evaluated system reaches; the \textbf{coefficient of variation}
$\sigma/\mu$ of performance across all configurations, a scale-invariant measure of relative
spread; and the \textbf{best-model baseline} $1 - p_{\text{best}}$, where $p_{\text{best}}$ is the
strongest single-model performance on the dataset.

\subsection{Domain characterisation}
Table~\ref{tab:complexity} gives the components. MedQA is a low-complexity domain ($D = 0.097$);
MedXpertQA and MedAgentsBench-hard are high-complexity and nearly indistinguishable
($D = 0.489$ and $0.478$).

\begin{table}[h]\centering\small
\caption{Domain complexity of the three medical benchmarks.}
\label{tab:complexity}
\begin{tabular}{@{}lrrrrr@{}}
\toprule
Benchmark & $1-p_{\max}$ & $\sigma/\mu$ & $1-p_{\text{best}}$ & $D$ & mean MAS gain \\
\midrule
MedXpertQA & 0.360 & 0.334 & 0.396 & \textbf{0.363} & +1.36 pp \\
MedAgentsBench-hard & 0.412 & 0.328 & 0.436 & \textbf{0.392} & -0.34 pp \\
MedQA (USMLE) & 0.036 & 0.108 & 0.036 & \textbf{0.060} & +0.72 pp \\
\bottomrule
\end{tabular}
\end{table}

\subsection{The critical threshold does not transfer}
\citet{kim2026scaling} identify a critical complexity threshold at $D \approx 0.40$: below it
multi-agent architectures yield net positive returns, above it coordination overhead consumes
resources otherwise spent on reasoning. Our data do not support a benchmark-level threshold of any
value. MedXpertQA ($D = 0.489$) and MedAgentsBench-hard ($D = 0.478$) differ by $0.011$ in
complexity yet by $3.5$ pp in mean multi-agent gain ($+3.09$ pp versus $-0.37$ pp) --- opposite
signs at effectively the same $D$. A benchmark-level statistic cannot separate them because the
quantity that does separate them varies \emph{within} each benchmark: on MedAgentsBench-hard the
per-tier gains run $-2.78$, $+1.02$, $+0.66$ pp as $\PSA$ moves from 20.4\% through 32.4\% to
44.4\%, entering the window between the first tier and the second. Domain complexity is the right
idea at the wrong granularity for medicine; the configuration-level baseline $\PSA$ is what
carries the signal.

\section{Datasets}
\label{app:data}

\textbf{MedXpertQA-Text}~\citep{zuo2025medxpertqa}: 2{,}450 expert-level items, ten options each,
tagged by 12 body systems and 3 task types. We draw a stratified sample interleaved across
(body system $\times$ task) so that any prefix is approximately balanced, then take the first 250.
\textbf{MedQA-USMLE}~\citep{jin2021medqa}: 1{,}273 four-option test items; randomly shuffled with a
fixed seed, first 250. \textbf{MedAgentsBench-hard}~\citep{medagentsbench}: the
contamination-screened hard split, 894 items after excluding the \texttt{medqa\_5options} subset,
stratified-interleaved across the ten source subsets, first 250. All samples are nested prefixes,
so a larger sample re-uses every call of a smaller one.

Item difficulty is additionally tagged easy/medium/hard by the pass rate of three mid-tier models
at $k{=}5$; on the MedXpertQA sample this yields 13/62/125 items.

\section{Implementation details}
\label{app:impl}

\subsection{Technical infrastructure}
All models are accessed through one OpenAI-compatible client with a content-addressed disk cache,
exponential backoff with rate-limit-aware retry, a per-call JSONL ledger recording tokens, latency
and cost, and a cross-process spend journal guarded by an OS file lock. Every LLM call is cached
on the full request (model, system, messages, temperature, seed, max tokens, effort, JSON mode,
tag), so a re-run is a replay.

\subsection{Agent configuration}
Single agents answer in one call. Independent panels deploy $N$ specialists with
synthesis-only aggregation. Centralized systems use one orchestrator over $N$ sub-agents with a
maximum of 2 orchestration rounds and one re-task per round. Decentralized systems run $N$ agents
through up to 3 discussion rounds, stopping early only on true unanimity. Hybrid systems combine a
generalist triage gate with the specialist panel and escalate to full discussion when the panel
has no majority. Agent temperature is 0.7 within panels, 0.3 for single-doctor baselines and for
orchestrator/moderator turns; reasoning models are called at \texttt{effort=low} with a minimum
completion budget of 1{,}600 tokens, since reasoning tokens are drawn from the same budget and a
smaller cap returns empty output.

\subsection{Prompt compilation}
Every agent receives the same rendered item (stem, then options as \texttt{(A) ...}) and the same
answer schema: JSON with \texttt{answer}, \texttt{confidence} (0--100) and \texttt{reason}
(at most 45 words). Only the system prompt differs, naming the specialty. Peer context in
communicating architectures is capped at a fixed 900-character budget, balanced-truncated across
peers and \emph{independent of $N$}, so a panel-size effect cannot be a context-length effect.
Peer ordering is deterministically rotated per (item, round) so that the first-listed speaker is
decorrelated from the top-ranked specialty.

\subsection{Evaluation methodology}
\textbf{Sample sizes.} 250 items per benchmark, 180 multi-agent configurations, 63{,}499 episodes.
\textbf{Controls.} All architectures share the item rendering, answer schema, role-prompt template
and decoding parameters; only coordination logic differs. Model weights are frozen. The primary
endpoint is exact match against the dataset's own gold option key --- there is no LLM judge in the
evaluation path. Episodes that fail for infrastructure reasons carry an error status and are
excluded from accuracy rather than scored as wrong, which would bias call-heavy architectures
downward. Cost accounting is nominal: cached calls are charged at list price, so reported
economics describe what a configuration would cost to run.

\subsection{Information gain computation}
Following \citet{kim2026scaling}, information gain is the Bayesian posterior variance reduction
\begin{equation}
\Delta\mathcal{I} = \tfrac12 \log \frac{\mathrm{Var}[Y \mid \mathbf{s}_{\text{pre}}]}
{\mathrm{Var}[Y \mid \mathbf{s}_{\text{post}}]}, \qquad
\mathrm{Var}[Y \mid \mathbf{s}] = \hat p(\mathbf{s})\,(1 - \hat p(\mathbf{s})),
\end{equation}
with $Y \in \{0,1\}$ the success indicator. They estimate $\hat p$ from $K{=}10$ sampled reasoning
traces at temperature 0.7. In multiple-choice consultation the panel already supplies such a
sample: $\hat p$ is the fraction of the round's opinions that name the gold option, Laplace
smoothed as $(k{+}1)/(n{+}2)$. Smoothing rather than clipping matters: a Centralized panel's final
round contains only the orchestrator, and clipping would pin $\hat p$ at $0.5$ and force
$\Delta\mathcal{I}$ negative by construction.
